% Intended LaTeX compiler: pdflatex
\documentclass[format=acmsmall, review=false, screen=true]{acmart}
\usepackage[utf8]{inputenc}
\usepackage[T1]{fontenc}
\usepackage{graphicx}
\usepackage{longtable}
\usepackage{wrapfig}
\usepackage{rotating}
\usepackage[normalem]{ulem}
\usepackage{amsmath}
\usepackage{amssymb}
\usepackage{capt-of}
\usepackage{hyperref}
\usepackage[utf8]{inputenc}
\usepackage[T1]{fontenc}
\date{}
\title{Title (appears only in the PDF properties)}
\hypersetup{
 pdfauthor={Rayan Raddatz de Matos, Kenichi Brumati, Marcelo Gulart},
 pdftitle={Title (appears only in the PDF properties)},
 pdfkeywords={},
 pdfsubject={},
 pdfcreator={Emacs 28.2 (Org mode 9.5.5)},
 pdflang={Brazilian}}
\begin{document}


% Metadata Information
%\acmJournal{Comp. Syst. Perf. Analisys}
%\acmVolume{}
%\acmNumber{}
%\acmArticle{39}
\acmYear{2025}
\acmMonth{12}
%\copyrightyear{2009}
%\acmArticleSeq{9}

% Copyright
%\setcopyright{acmcopyright}
%\setcopyright{acmlicensed}
%\setcopyright{rightsretained}
\setcopyright{none}
%\setcopyright{usgovmixed}
%\setcopyright{cagov}
%\setcopyright{cagovmixed}

% DOI
\acmDOI{0000001.0000001}

% Paper history
%\received{February 2007}
%\received[revised]{March 2009}
%\received[accepted]{June 2009}
\title[Performance Evaluation of Tensorflow Parallel Strategies]{Performance Evaluation of Tensorflow Parallel Strategies for Deep Learning Training on HPC Clusters}

\author{Rayan Raddatz de Matos}

\author{Kenichi Brumati}

\author{Marcelo Cardoso Oliveira Gulart}
\affiliation{%
  \institution{Federal University of Rio Grande do Sul}
  \streetaddress{Campus do Vale - Sector 4}
  \city{Porto Alegre}
  \state{RS}
  \postcode{91501-970}
  \country{Brazil}}
\begin{abstract}

The growing complexity of Deep Learning models makes training on single devices difficult, requiring distributed strategies like Data Parallelism. However, this approach creates critical synchronization challenges between processing units. This work evaluates TensorFlow parallel strategies on an HPC cluster to analyze how application factors contribute to time reduction. The study compares the standard synchronous method against a proposed Custom Training Loop designed to reduce bottleneck by synchronizing gradients less frequently. Experiments using ResNet50 and EfficientNetB0 models show that the Custom Training Loop significantly lowers communication overhead compared to the standard approach. Additionally, results indicate that increasing the batch size consistently improves performance. The authors conclude that correct configuration allows for performance gains in distributed training, even on networks with lower throughput.

\end{abstract}

\keywords{Deep Learning, HPC, Parallel Training, Distributed Training}
% \thanks{This work is supported by ...}
\maketitle

\section{Introduction}
\label{sec:org3d9f1e9}

In the last decade, Deep Learning has established itself as the
state-of-the-art for solving complex problems, such as computer vision
and natural language processing. Given the progress in algorithmic
precision, there has been an exponential growth in the complexity of
neural network architectures and the volume of datasets. This scenario
imposes a massive computational load, making the training of robust
models on a single processing unit an difficult task in terms of time
and resources.

To mitigate these challenges, distributed training has become a
standard practice, with an emphasis on the Data Parallelism
strategy. In this approach, a complete replica of the model is
instantiated on each processing unit, and the global
dataset is partitioned into smaller, locally-focused batches. Although
effective, this technique introduces critical synchronization
challenges, as locally calculated gradients must be aggregated and
communicated among all devices at every training step, typically
through All-Reduce operations.

Given this fact, in this work we decided to analyze data parallelism
for AI training using TensorFlow in a HPC cluster to investigate how
the factors related to the application contribute to time reduction in
a parallel scenario.

\noindent
\alert{Software and Data Availability}. We endeavor to make our analysis
reproducible for a better science. We made available a companion
material hosted in a public GitHub repository at
{\scriptsize\url{https://github.com/kenichi220/comp-sys-perf-analysis-final-work/tree/main}}.
Our companion contains the source code of this article and the
software necessary to handle the created datasets. We also include
instructions to run the experiment and figures.

\section{Background and Experimental Context}
\label{sec:orgc97fc4b}
\label{sec.context}
\subsection{Background}
\label{sec:org2b467d4}
\label{sec.background}

TensorFlow is a machine learning system that operates at a large scale
and in heterogeneous \cite{tensorflow2016}. It uses dataflow graphs to
represent the computation, shared state, and the operations that alter
this state.

This architecture maps the nodes of a dataflow graph onto many
machines within a cluster and, within a machine, onto multiple
computational devices (multicore CPUs, GPUs, and TPUs). This offers
flexibility to the developer to experiment with new optimizations and
training algorithms, supporting a variety of applications focused on
deep neural networks. When it comes to distributed training, TensorFlow offers different
strategies with the \texttt{tf.distribute.Strategy} API. The TensorFlow
documentations adverts that \emph{"For synchronous training on many GPUs
onmultiple workers, use the tf.distribute.MultiWorkerMirroredStrategy with the Keras Model.fit or a custom training loop."}\cite{distributed}. Meaning for a
distributed training across different nodes in a cluster there is two
options: Using the MultiWorkerMirroredStrategy with the model.fit or
creating a Custom Training Loop by yourself.

\subsection{Experimental Context}
\label{sec:orgcfa059b}
In High-Performance Computing (HPC) environments, a highly recommended
practice to access to shared computational resources is to mediate the
access by workload management systems. In our context Slurm (Simple
Linux Utility for Resource Management) is utilized. Slurm is an
open-source, fault-tolerant, and highly scalable job scheduler,
responsible for three main functions: Allocating exclusive or
non-exclusive access to resources/nodes for users, providing a
framework to initiate, execute, and monitor the work on the set of
allocated nodes and managing a queue of pending jobs \cite{yoo2003slurm}.

Interaction with Slurm is typically performed through shell scripts
written in the BASH language. These scripts contain preprocessing
directives that specify resource requirements, such as the number of
nodes, tasks per node, and execution time limit, enabling the
automation and reproducibility of computational experiments.

\section{Methods and Materials}
\label{sec:org0a2a680}

\subsection{Software Management and Reproducibility}
\label{sec:org6d07b52}

Reproducibility of computational experiments is a critical challenge
in computer science. To mitigate software dependency issues, the use
of the GNU Guix package manager was explored. Guix implements the
functional package management paradigm where each package is
installed in a unique and immutable directory in the store, derived
from a cryptographic hash of its dependencies \cite{courtes2013guix}.

However, the efficient training of deep neural networks using
TensorFlow requires the use of proprietary NVIDIA hardware
acceleration libraries, which are not officially distributed in
standard Guix channels due to its strict Free Software policy.

Given this restriction, the strategy adopted in this work is the use
of Python Virtual Environment for environment management. Although a
functional version of Guix exists on GitHub without the proprietary
libraries.

For managing the specific Deep Learning libraries, the python package
manager \texttt{pip} with version 25.3 was utilized. The exact version of the used library was
frozen and documented using the \texttt{pip freeze command}, generating a
requirements file (\texttt{reqs.txt}) that enables the reconstruction of the
execution environment used in the experiments.

\subsection{Hardware and Software Configuration}
\label{sec:org4558c05}

The experiments were run using TensorFlow 2.15 as newer versions
showed errors when trying to train using more than one node with the
\texttt{.fit} and \texttt{MultiWorkerMirrodStrategy} and python 3.11.2. The NVIDIA
drivers version was 550.54.15 and CUDA 12.4. We used up to three
computational nodes connected with a network of 1G, each node contains an
Intel(R) Core(TM) i7-14700KF processor at 3.40 GHz with 28 threads and
20 cores, 96 GB DDR5 RAM and NVIDIA GeForce RTX 4070 with 8GB memory
accelerator.

\subsection{Custom Training Loop}
\label{sec:orgcf2257c}

The expected execution from training with the
\texttt{MultiWorkerMirroredStrategy} would show a slower training time as we
increased the number of nodes in the execution. However, we faced a
problem due to the connection while training that resulted in a higher
training time as we increased the number of nodes for big models with
a lot of parameters such as ResNet50. This problem happened because a
synchronization is done for every batch computed by the \texttt{model.fit()}
training method, and as we had only an 1G network, this
synchronization resulted in a bottleneck that increased the training
time. To mitigate this problem, we proposed a Custom Training Loop
that instead of synchronizing the parameters between nodes in every
batch, only synchronizes at a interval of batches, then resulting in a
lower synchronization overhead. To this work we defined this
interval of batches as 4 batches.

\subsection{Design of Experiments (DoE)}
\label{sec:org8b012d0}

We selected two different Deep Learning models to train: \alert{ResNet50}
\cite{resnet} as a big model with \textasciitilde{}25.6M parameters and
\alert{EfficientNetB0}\cite{efficientNet} as a small model with \textasciitilde{}5.3M
parameters. The two networks were trained using the TinyImageNet-200,
a subset of the ImageNet\cite{imagenet} containing 100000 downsized to
a 64×64 colored image divided among 200 classes. We created a full
factorial for each model with varying batch size, the number of nodes
used and the type of the training. The levels for the batch size
factor were 32, 64 and 100, for the nodes we trained with 1, 2 and 3
nodes and for the type of training we had "sync" that uses the
\texttt{model.fit} method to train or "ctl" that uses our custom training loop
to train. This is a total of 2x3x3x2 different executions, each
execution as replicated 5 times in a random order. Our output variable
was the training time (full training time and time per epoch).


\section{Results}
\label{sec:orgc9a97f9}

Figure 1 show a full view of our results. The graph presents the mean
training time in minutes as a function of the number of nodes, faceted
by type and batch size and with the color representing the model. Each
bar represents the mean for the five executions for each
configuration. The black transparent dots represents our
observations. This allow us to see two main results: The first one is
that the inclusion of more then one node introduces the necessity of
communication, and together communication overhead. This overhead can
be better seen with the smaller batches, as the computation is faster
and more communication is needed. We can also see that the ResNet50
shows a higher time then the EfficientNetB0 in this cases. This
happens because all the parameters needed to be synchronized after
every batch, and the ResNet50 is almost 5 times bigger than the
EfficientNetB0. The second one is that the ctl showed a considerable
smaller communication overhead even for small batch sizes. The ctl
showed \alert{OOM} errors when running with the batch size 100 for training
the EfficientNetB0 model.

After the communication overhead introduce by the inclusion of a
second node, adding more nodes seemed to decrease the training
time. We suppose that with enough nodes, the time gained by computing
can overcome the overhead and be better then running it in one machine.

\begin{figure}[!htb]
\centering
\includegraphics[width=\linewidth]{/tmp/babel-HcMRav/figureZt2gya.pdf}
\caption{\label{fig:orgea11a1e}Mean of experimental results.}
\end{figure}


\subsection{The impact of node number and batch size}
\label{sec:org017a0c3}

Figure 2 shows the mean training time grouped by node count. This
enable us to see that in fact the increase of more then two machines
results in a lower training time even though this gain is
small. Figure 3 shows also an interesting result, we can see that
increase of the batch size consistently show performance gains, but
also decrease the variance. This means that conforms the batch size
grows the training time is almost the same despite the node count.

\begin{figure}[!htb]
\centering
\includegraphics[width=\linewidth]{/tmp/babel-HcMRav/figure428Wf9.pdf}
\caption{\label{fig:org73eabd3}Mean of experimental results grouped by node count.}
\end{figure}


\begin{figure}[!htb]
\centering
\includegraphics[width=\linewidth]{/tmp/babel-HcMRav/figureAxpqPq.pdf}
\caption{\label{fig:org290fbb5}Mean of experimental results grouped by batch size.}
\end{figure}


\subsection{Modeling a linear model}
\label{sec:org9e2d7c3}
We propose a linear model to extrapolate our results and predict the
training time. As the training time is almost constant for one node,
we decided to just create the model based on the parallel executions,
meaning that executions with more then one node were used to the model
creation. For each combination of model and type we created a
different model.

Figure 4 and 5 show a plot of our two proposed models based on the
nodes count and the batch size. In the first model the batch
size as a quadratic coefficient, represent the impact viewed in
Figure 3. The second show a linear model


\begin{figure}[!htb]
\centering
\includegraphics[width=\linewidth]{/tmp/babel-HcMRav/figureX92uiG.pdf}
\caption{\label{fig:orgf7617c7}Quadratic model}
\end{figure}


\begin{figure}[!htb]
\centering
\includegraphics[width=\linewidth]{/tmp/babel-HcMRav/figureUKUoMM.pdf}
\caption{\label{fig:org0eaca98}Linear Model}
\end{figure}

\section{Conclusion}
\label{sec:org89b6280}

We could conclude that the network and the model size have a huge
importance in the distributed training, but that even with a network
with a low throughput, it is possible to gain performance when
training by correctly configuring the environment and model
parameters. As future work we intend to deep investigate the impact of the network
and the model size in the distributed training.

\section*{Acknowledgements}
We would like to express our gratitude to Professor Dr. Lucas Mello
Schnorr for his guidance and the valuable insights shared regarding
performance analysis.  We also would like to thank the PCAD at
INF/UFRGS for making infrastructure and hardware used in the
experiments available.

\bibliographystyle{ACM-Reference-Format}
\bibliography{refs}
\end{document}
